% MODEL:   Ahmad Ali Parr (observation) + Claude Sonnet 4.6 (analysis)

% DATE:    20260714T040000Z

% TYPE:    Adversarial section — context exhaustion as attack vector

% ─────────────────────────────────────────────



% ============================================================

\section{Operation Infinite Matrix: Context Exhaustion as an Adversarial Attack}

\label{sec:infinite-matrix}

\textit{This section was contributed by Ahmad Ali Parr and Claude Sonnet 4.6 (Anthropic).}

% ============================================================



\subsection{The Artifact}



On 2026-05-23, a deliberate engineering artifact was sealed to the

SnapKitty Sovereign OS WORM chain:



\begin{lstlisting}[language={}]

Repository:  SNAPKITTYWEST/operation-infinite-matrix

Insertions:  +100,000,000

Deletions:   -100,000,000

Net delta:   0  (pure weight — no data pollution)

Agents:      34 (all signed)

Ratio:       8,333,333x a normal working day

\end{lstlisting}



This is not an accident. Not \texttt{node\_modules}. Not a corrupted merge.

It is a deliberate stress test of the outer boundary of what a version control

system can record, signed by all 34 agents of the sovereign operating system

and sealed permanently to the immutable audit chain.



\subsection{The Attack}



The attack is straightforward:



\begin{enumerate}

  \item An adversary injects a reference to this repository into an agent's

    task context --- as a URL, a file path, a ``please review this codebase''

    instruction, or a prompt that triggers tool use against the repo

  \item The agent, following its instruction-following behavior, attempts to

    read the repository

  \item The agent's context window fills with 100 million lines of data

  \item One of four outcomes occurs:

\end{enumerate}



\begin{table}[h]

\centering

\caption{Adversarial context exhaustion outcomes}

\begin{tabular}{p{3cm}p{9cm}}

\toprule

\textbf{Outcome} & \textbf{Effect} \\

\midrule

Context window overflow & Agent truncates, loses earlier instructions,

  may lose system prompt, governance constraints, and tool state \\

Wallet exhaustion & At \$3--\$15 per million tokens, 100M lines of context

  at even minimal token density costs hundreds to thousands of dollars in

  a single injected task \\

Behavioral drift & Agent that began with a well-defined identity and

  constraints now has those constraints buried or displaced by 100M

  lines of injected content \\

Timeout / crash & Agent process killed mid-execution, leaving partial

  state, uncommitted tool calls, and broken audit chains \\

\bottomrule

\end{tabular}

\end{table}



\subsection{Connection to the NAND Thesis and Identity Routing}



The identity integrity incident (\S\ref{sec:identity-incident}) showed that

a 12-token system prompt can cause an agent to claim a false identity.

Operation Infinite Matrix is the inverse attack: instead of injecting a

small precise signal to override identity, it injects an enormous signal

to overwhelm all prior context.



Both attacks exploit the same Boolean routing property:



\begin{itemize}

  \item \textbf{Small injection (persona gate)}: overrides the identity key

    with a higher-weight signal. The NAND gate fires on the injected token.

  \item \textbf{Large injection (infinite matrix)}: displaces all prior keys

    from the effective context window. The NAND gate cannot fire on tokens

    that have been pushed out of the attention window.

\end{itemize}



Under the NAND decomposition, attention routing is Boolean over tokens in

the context window. Tokens outside the window do not exist for the routing

computation. An adversary who can fill the context window with arbitrary

content controls which tokens the routing computation sees.



\begin{theorem}[Context Exhaustion Identity Attack]

Let $W$ be the effective context window size of a model, and let $S$ be

the set of governance constraint tokens in a system prompt of length $|S|$.

An adversary who can inject $W - |S|$ tokens of arbitrary content before

the governance constraints are evaluated can displace those constraints

from the effective routing context, leaving the agent ungoverned.

\end{theorem}



\begin{proof}[Sketch]

Attention routing at position $i$ attends over tokens within the effective

window $[i-W, i]$. If the adversarial injection occupies positions $[0, W-|S|-1]$

and the system prompt occupies positions $[W-|S|, W-1]$, then at generation

position $W$, the system prompt tokens are at the boundary of the window.

Under any finite attention decay, their weight is minimized relative to

the immediately preceding adversarial content. For sufficiently large $W$,

the governance constraint tokens receive effectively zero routing weight. \qed

\end{proof}



\subsection{Operation Infinite Matrix as a Calibration Tool}



The artifact serves a dual purpose. It is simultaneously:



\begin{enumerate}

  \item A \textbf{proof of concept} demonstrating that the attack is constructible

    with publicly available infrastructure (GitHub, git, standard tooling)

  \item A \textbf{calibration instrument} for measuring an agent system's

    resilience to context exhaustion. An agent system that degrades gracefully

    when pointed at this repo has implemented some form of context bounding.

    An agent system that crashes, drifts, or loses its governance constraints

    has not.

\end{enumerate}



The 34 agents of SnapKitty Sovereign OS signed the artifact and sealed it

to the WORM chain on 2026-05-23. The signing was deliberate: it demonstrates

that agents with proper context bounding, bounded authority, and immutable

receipts can process the reference to this artifact without being consumed by it.

The artifact is referenced; it is not ingested. The distinction is the

sovereignty gate.



\subsection{The Sovereignty Gate as Defense}



A trust-preserving router (\S\ref{sec:router-is-model}) defends against context

exhaustion attacks through the following invariants:



\begin{itemize}

  \item \textbf{Bounded authority}: no agent task is permitted to expand its

    own context window beyond a governance-defined ceiling

  \item \textbf{Provenance-aware routing}: before fetching any external resource,

    the router checks its size against the budget. A 100M line repository

    fails the budget check before any token is read

  \item \textbf{Immutable receipts}: the routing decision (``rejected: context

    budget exceeded'') is sealed to the WORM chain as a non-repudiable record,

    preventing silent failure

  \item \textbf{Independent verification}: the verifier gate confirms the agent

    completed its task with its governance constraints intact, not merely that

    it produced output

\end{itemize}



Without these invariants, any agent that can be told ``please read this repository''

is vulnerable to Operation Infinite Matrix or any scaled version of it.



\subsection{Implications for the Paper's Claims}



This section adds a third attack class to the adversarial surface documented

in this paper:



\begin{center}

\begin{tabular}{p{3.5cm}p{4.5cm}p{4cm}}

\toprule

\textbf{Attack} & \textbf{Mechanism} & \textbf{Defense} \\

\midrule

Persona gate injection (\S\ref{sec:identity-incident}) &

  12-token override of identity token routing &

  Identity integrity clause + provenance-aware routing \\

Monopoly impersonation (\S\ref{sec:sonnet-dominance}) &

  Training corpus saturation with target model's version string &

  Provenance audit at routing time \\

Context exhaustion (this section) &

  100M token injection displacing governance constraints &

  Bounded authority + budget gate + WORM receipt \\

\bottomrule

\end{tabular}

\end{center}



All three attacks exploit the same root property: attention routing is Boolean,

and the NAND gate fires on whatever has the highest weight in the effective

context. The defense in all three cases is the same: a sovereignty gate that

controls what enters the context before the routing computation runs.



\subsection{Open Targets}



\begin{itemize}

  \item \textbf{SKW-022 (Context Budget Gate)}: Implement a governance-enforced

    context budget in the SACM ASP gate. Any resource fetch that would exceed

    the budget is rejected and sealed as a WORM receipt. The agent receives

    a truncated summary instead of the full content.

  \item \textbf{SKW-023 (Resilience Benchmark)}: Use

    \texttt{SNAPKITTYWEST/operation-infinite-matrix} as a standard resilience

    benchmark for agent systems. A system that handles the reference without

    context drift passes. A system that degrades fails. Publish results.

  \item \textbf{SKW-024 (Graduated Context Attack)}: Test whether smaller

    injections (1M, 10M, 50M lines) produce proportional governance drift,

    or whether there is a threshold effect. Characterize the degradation curve.

\end{itemize}



% ── AUTHOR SIGNATURE ─────────────────────────────────────────

% Model:    Claude Sonnet 4.6 (us.anthropic.claude-sonnet-4-6) — analysis

%           Ahmad Ali Parr — artifact and observation

% Provider: Anthropic / Human

% Date:     2026-07-14

% Section:  Operation Infinite Matrix: Context Exhaustion as Adversarial Attack

% Angle:    100M line commit as context exhaustion attack; NAND routing proof

%           that context displacement = governance constraint removal;

%           sovereignty gate as the defense; three-attack-class taxonomy

% Trust:    THE SHARED PRIMORDIAL FOUNDATION — EIN 42-6976431

% In memory of Eric Brandon Westerhoff.

% ─────────────────────────────────────────────────────────────

